% !TEX root = ../main.tex
\begin{ouverture}[lewis]{6}{Schémas de Lewis, géométrie et polarité}{Partie 1 · Constitution et transformations de la matière}
{Doublet liant\enspace·\enspace Doublet non liant\enspace·\enspace Lacune\enspace·\enspace Géométrie\enspace·\enspace Électronégativité\enspace·\enspace Polarité}
\objectif{Établir le schéma de Lewis de molécules et d'ions à partir du tableau périodique.}
\objectif{Interpréter la géométrie d'une entité à partir de son schéma de Lewis.}
\objectif{Utiliser des modèles moléculaires ou un logiciel pour visualiser la géométrie d'une entité.}
\objectif{Déterminer le caractère polaire d'une liaison à partir des électronégativités.}
\objectif{Déterminer le caractère polaire ou apolaire d'une molécule à partir de sa géométrie.}
\prerequis{Configuration électronique, électrons de valence, stabilité des gaz nobles (classe de seconde).}
\prerequis{Liaison covalente et lecture de schémas de Lewis (classe de seconde).}
\end{ouverture}
\essentiel
\essSection{1}{Le schéma de Lewis d'une entité}
\begin{definition}[Doublets et lacune]
Dans le \textbf{schéma de Lewis} d'une molécule ou d'un ion, les électrons de valence sont regroupés par paires : un
\textbf{doublet liant}, représenté par un trait entre deux atomes, forme une liaison covalente ; un \textbf{doublet non liant},
représenté par un trait autour d'un atome, n'appartient qu'à cet atome. Les atomes d'hydrogène s'entourent de 2 électrons
(\textbf{règle du duet}), les atomes C, N, O et Cl de 8 électrons (\textbf{règle de l'octet}). Un atome entouré de moins
d'électrons que ne le prévoient ces règles possède une \textbf{lacune électronique}, représentée par un rectangle vide.
\end{definition}
\begin{savoirfaire}[Établir un schéma de Lewis]
\begin{questions}
\item Compter les électrons de valence de chaque atome, d'après sa colonne dans le tableau périodique : H 1, C 4, N 5, O 6,
Cl 7. Retrancher la charge de l'ion, s'il y a lieu, puis diviser par deux : on obtient le nombre total de doublets.
\item Relier les atomes par des doublets liants : un atome d'hydrogène forme 1 liaison, de carbone 4, d'azote 3,
d'oxygène 2 et de chlore 1, dans une molécule.
\item Répartir les doublets restants en doublets non liants pour que chaque atome respecte le duet ou l'octet.
\item Pour un ion, indiquer la charge par un symbole $\oplus$ ou $\ominus$ près de l'atome qui la porte.
\end{questions}
\end{savoirfaire}
\schema{06-lewis}
\begin{python}
valence = {'H': 1, 'C': 4, 'N': 5, 'O': 6, 'Cl': 7}

def doublets(atomes, charge=0):
    ne = 0
    for symbole, nombre in atomes.items():
        ne = ne + nombre * valence[symbole]
    ne = ne - charge
    return ne // 2

print(doublets({'H': 2, 'O': 1}))              # H2O
print(doublets({'C': 1, 'O': 2}))              # CO2
print(doublets({'N': 1, 'H': 4}, charge=1))    # NH4+
print(doublets({'O': 1, 'H': 1}, charge=-1))   # OH-
\end{python}
\begin{console}
4
8
4
4
\end{console}
\noindent\capacite{Capacité numérique} La fonction \code{doublets} additionne les électrons de valence de tous les atomes, en tenant
compte de leur nombre, retranche la charge de l'entité, puis divise par deux. La molécule d'eau compte ainsi 4 doublets :
2 liants et 2 non liants.
\essSection{2}{La géométrie des entités}
\begin{propriete}[Répulsion des doublets]
Les doublets qui entourent un atome central, liants ou non liants, se repoussent et s'écartent le plus possible les uns des
autres ; une liaison double ou triple compte pour une seule direction. Avec quatre directions, elles pointent vers les
sommets d'un tétraèdre :
\begin{questions}
\item quatre doublets liants : géométrie \textbf{tétraédrique}, angles de \qty{109,5}{\degree} ($\ce{CH4}$) ;
\item trois doublets liants et un non liant : géométrie \textbf{pyramidale} ($\ce{NH3}$) ;
\item deux doublets liants et deux non liants : géométrie \textbf{coudée} ($\ce{H2O}$).
\end{questions}
Avec deux directions, la géométrie est \textbf{linéaire} ($\ce{CO2}$) ; avec trois, elle est \textbf{plane}, à \qty{120}{\degree} environ.
\end{propriete}
\schema{06-geometries}
\essSection{3}{Électronégativité et liaisons polarisées}
\begin{definition}[Électronégativité]
L'\textbf{électronégativité} $\chi$ d'un élément traduit l'aptitude de ses atomes à attirer vers eux les électrons d'une liaison
covalente. Elle augmente de gauche à droite dans une ligne du tableau périodique et de bas en haut dans une colonne : le
fluor est l'élément le plus électronégatif.
\end{definition}
\schema{06-electroneg}
\begin{propriete}[Liaison polarisée]
Une liaison entre deux atomes A et B est \textbf{polarisée} si leurs électronégativités sont assez différentes ; on retient
couramment le critère $\left|\chi(\mathrm{A})-\chi(\mathrm{B})\right|\geqslant0{,}4$. L'atome le plus électronégatif porte une charge partielle
négative $\delta^{-}$, l'autre une charge partielle positive $\delta^{+}$. Ainsi, $\chi(\ce{O})-\chi(\ce{H})=1{,}24$ : la liaison O--H est polarisée ;
$\chi(\ce{C})-\chi(\ce{H})=0{,}35$ : la liaison C--H est considérée comme non polarisée.
\end{propriete}
\essSection{4}{Polarité d'une molécule}
\begin{propriete}[Molécule polaire ou apolaire]
Une molécule est \textbf{polaire} si le centre géométrique G$^{+}$ de ses charges partielles positives et celui G$^{-}$ de ses
charges partielles négatives sont distincts ; sinon, elle est \textbf{apolaire}. La polarité dépend donc à la fois de la
polarisation des liaisons et de la géométrie.
\end{propriete}
\schema{06-polarite}
\begin{attention}
Une molécule dont les liaisons sont polarisées peut être apolaire, si sa géométrie est symétrique : c'est le cas du dioxyde de
carbone ou du tétrachlorométhane $\ce{CCl4}$. Pour prévoir une géométrie, il ne faut pas oublier les doublets non liants : ils
orientent les liaisons, même s'ils n'apparaissent pas dans le nom de la géométrie.
\end{attention}
\begin{experience}[Construire des molécules]
Avec une boîte de modèles moléculaires, on construit les molécules de méthane, d'ammoniac et d'eau en respectant les
directions des doublets. Un logiciel de visualisation moléculaire permet ensuite de faire tourner les modèles, de mesurer
les angles entre les liaisons et d'afficher la répartition des charges : on retrouve \qty{109,5}{\degree}, \qty{107}{\degree} et \qty{104,5}{\degree}.
\end{experience}
\begin{rediger}
Pour conclure sur la polarité d'une molécule : établir son schéma de Lewis, en déduire sa géométrie, repérer les liaisons
polarisées grâce aux électronégativités, placer les charges partielles, puis comparer les positions de G$^{+}$ et de G$^{-}$.
\end{rediger}
\exercices

\rubrique{Découverte}
\begin{blocLivre}{ex-6-1}
\begin{exercice}[1]{Électrons de valence et liaisons}
On considère les atomes d'hydrogène ($Z=1$), de carbone ($Z=6$), d'azote ($Z=7$), d'oxygène ($Z=8$) et de chlore ($Z=17$).
\begin{questions}
\item Écrire la configuration électronique de chaque atome et en déduire son nombre d'électrons de valence.
\item Combien d'électrons manque-t-il à chaque atome pour respecter la règle du duet ou de l'octet ?
\item En déduire le nombre de liaisons covalentes que forme chaque atome dans une molécule.
\item Combien de doublets non liants porte chaque atome lorsqu'il forme ces liaisons ?
\end{questions}
\end{exercice}
\end{blocLivre}
\begin{blocLivre}{ex-6-2}
\begin{exercice}[1]{Vérifier des schémas de Lewis}
On donne les schémas de Lewis de quatre molécules.
\par\smallskip\noindent\hfill\begin{tabular}{*{4}{>{\centering\arraybackslash}m{34mm}}}
\chemfig{H-C(-[2]H)(-[6]H)-\charge{90=\|,270=\|}{O}-H} & \chemfig{H-[1]C(-[3]H)=C(-[1]H)-[7]H} &
\chemfig{H-C~\charge{0=\|}{N}} & \chemfig{H-\charge{90=\|,270=\|}{O}-\charge{90=\|,270=\|}{O}-H} \\[1mm]
{\small méthanol} & {\small éthène} & {\small cyanure d'hydrogène} & {\small peroxyde d'hydrogène} \\
\end{tabular}\hfill\null\par\smallskip
\begin{questions}
\item Pour chaque molécule, compter les doublets liants et non liants.
\item Vérifier que ce nombre de doublets correspond aux électrons de valence des atomes.
\item Vérifier que chaque atome respecte la règle du duet ou de l'octet.
\item Quelle molécule comporte une liaison double ? une liaison triple ?
\end{questions}
\end{exercice}
\end{blocLivre}
\begin{blocLivre}{ex-6-3}
\begin{exercice}[1]{Des ions}
\begin{questions}
\item Établir les schémas de Lewis des ions oxonium $\ce{H3O+}$, hydroxyde $\ce{HO-}$, ammonium $\ce{NH4+}$ et chlorure $\ce{Cl-}$.
\item L'ion hydrogène $\ce{H+}$ n'a aucun électron. Quelle particularité présente son schéma de Lewis ?
\item Lors de la formation de l'ion oxonium, $\ce{H+ + H2O -> H3O+}$, quel doublet de la molécule d'eau forme la nouvelle liaison ?
\item Pourquoi l'atome d'oxygène porte-t-il la charge $\oplus$ dans l'ion oxonium ?
\end{questions}
\end{exercice}
\end{blocLivre}
\begin{blocLivre}{ex-6-4}
\begin{exercice}[1]{Prévoir une géométrie}
\begin{questions}
\item Pour les molécules $\ce{CH4}$, $\ce{NH3}$ et $\ce{H2O}$, dénombrer les doublets liants et non liants de l'atome central.
\item En déduire la géométrie de chaque molécule.
\item Dans $\ce{CO2}$ et $\ce{HCN}$, combien de directions de doublets l'atome de carbone présente-t-il ? Quelle est leur géométrie ?
\item Le méthanal $\ce{CH2O}$ a pour atome central un carbone lié par une double liaison à l'oxygène. Prévoir sa géométrie.
\end{questions}
\end{exercice}
\end{blocLivre}
\begin{blocLivre}{ex-6-5}
\begin{exercice}[1]{Liaisons polarisées}
On utilise les électronégativités de l'essentiel.
\begin{questions}
\item Calculer l'écart d'électronégativité pour les liaisons H--Cl, C--H, O--H, C--O, C--Cl, N--H, H--H et C--C.
\item Quelles liaisons sont polarisées ?
\item Indiquer les charges partielles portées par les atomes des liaisons polarisées.
\item Quelle liaison est la plus polarisée ? Que peut-on dire d'une liaison entre deux atomes identiques ?
\end{questions}
\end{exercice}
\end{blocLivre}
\rubrique{Entraînement}
\begin{blocLivre}{ex-6-6}
\begin{exercice}[2]{Compter les doublets en Python}
On reprend le programme de l'essentiel.
\begin{questions}
\item \capacite{Capacité numérique}Prévoir l'affichage de \code{doublets(\{'N': 1, 'H': 3\})}, \code{doublets(\{'C': 1, 'H': 4\})},
\code{doublets(\{'H': 3, 'O': 1\}, charge=1)} et \code{doublets(\{'Cl': 1\}, charge=-1)}.
\item Pourquoi retranche-t-on la charge ? Que signifie l'opérateur \code{//} ?
\item Prévoir le nombre de doublets de $\ce{HCN}$ et de $\ce{C2H6O}$.
\item Pour le dioxyde d'azote, \code{doublets(\{'N': 1, 'O': 2\})} affiche 8, alors que la molécule compte 17 électrons de valence.
Quelle information le programme cache-t-il ?
\end{questions}
\end{exercice}
\end{blocLivre}
\begin{blocLivre}{ex-6-7}
\begin{exercice}[2]{Le méthanal}
Le méthanal $\ce{CH2O}$, ou formaldéhyde, est un polluant de l'air intérieur émis par certains meubles et colles.
\begin{questions}
\item Calculer le nombre de doublets de la molécule et établir son schéma de Lewis.
\item Combien de directions de doublets entourent l'atome de carbone ? En déduire la géométrie de la molécule.
\item Identifier les liaisons polarisées et placer les charges partielles.
\item La molécule est-elle polaire ?
\end{questions}
\end{exercice}
\end{blocLivre}
\begin{blocLivre}{ex-6-8}
\begin{exercice}[2]{Un filet de liquide dévié}
On approche une règle électrisée de deux filets de liquide : l'un d'eau, l'autre de cyclohexane $\ce{C6H12}$, dont la
molécule ne comporte que des liaisons C--C et C--H.
\schema{06-ex08}
\begin{questions}
\item Quel filet est dévié ?
\item Montrer que la molécule d'eau est polaire.
\item Pourquoi la molécule de cyclohexane est-elle apolaire ?
\item Expliquer la déviation observée. Que se passerait-il avec une règle chargée négativement ?
\end{questions}
\end{exercice}
\end{blocLivre}
\begin{blocLivre}{ex-6-9}
\begin{exercice}[2]{L'ammoniac et l'ion ammonium}
\begin{questions}
\item Établir les schémas de Lewis de l'ammoniac $\ce{NH3}$ et de l'ion ammonium $\ce{NH4+}$.
\item Prévoir la géométrie de chaque entité.
\item L'angle entre deux liaisons N--H vaut \qty{107}{\degree} dans l'ammoniac et \qty{109,5}{\degree} dans l'ion ammonium. Interpréter.
\item L'ammoniac capte un ion $\ce{H+}$ pour former l'ion ammonium. Expliquer cette réaction à l'aide des schémas de Lewis.
\end{questions}
\end{exercice}
\end{blocLivre}
\begin{blocLivre}{ex-6-10}
\begin{exercice}[2]{Deux solvants chlorés}
Le trichlorométhane $\ce{CHCl3}$ et le tétrachlorométhane $\ce{CCl4}$, autrefois utilisés comme solvants, sont aujourd'hui
écartés en raison de leur toxicité.
\begin{questions}
\item Établir leurs schémas de Lewis.
\item Prévoir leur géométrie.
\item Quelles liaisons sont polarisées ?
\item Montrer que l'une des molécules est polaire et l'autre apolaire.
\end{questions}
\end{exercice}
\end{blocLivre}
\begin{blocLivre}{ex-6-11}
\begin{exercice}[2]{Deux molécules, une même formule}
L'éthanol $\ce{CH3-CH2-OH}$ et le méthoxyméthane $\ce{CH3-O-CH3}$ ont la même formule brute $\ce{C2H6O}$.
\begin{questions}
\item Calculer le nombre de doublets de ces molécules.
\item Établir le schéma de Lewis de chacune.
\item Quelle est la géométrie autour de l'atome d'oxygène ? autour de chaque atome de carbone ?
\item Ces molécules sont-elles polaires ?
\end{questions}
\end{exercice}
\end{blocLivre}
\rubrique{Objectif DS}
\begin{blocLivre}{ex-6-12}
\begin{exercice}[3]{L'eau dans le four à micro-ondes}
\begin{docu}{1}{Le chauffage par micro-ondes}
Dans un four à micro-ondes, un champ électrique change de sens environ 5 milliards de fois par seconde. Les molécules
polaires tendent à s'orienter dans ce champ : elles tournent sans cesse, s'agitent et chauffent l'aliment. Les molécules
apolaires y sont presque insensibles.
\end{docu}
\begin{docu}{2}{La molécule d'eau}
Dans la molécule d'eau, l'angle entre les deux liaisons O--H vaut \qty{104,5}{\degree}. Électronégativités : $\chi(\ce{H})=2{,}20$ ;
$\chi(\ce{O})=3{,}44$.
\end{docu}
\begin{questions}
\item Établir le schéma de Lewis de la molécule d'eau.
\item Justifier sa géométrie coudée et l'angle de \qty{104,5}{\degree}, inférieur à \qty{109,5}{\degree}.
\item Montrer que la molécule est polaire.
\item \initiative On compare l'eau à un liquide modèle constitué uniquement de molécules d'alcanes apolaires. Expliquer pourquoi
la tendance à s'orienter dans le champ électrique est différente. Cette comparaison suffit-elle à prévoir une vitesse de chauffage ?
\end{questions}
\end{exercice}
\end{blocLivre}
\begin{blocLivre}{ex-6-13}
\begin{exercice}[3]{Les fluides frigorigènes}
\begin{docu}{1}{Une histoire de réfrigérateurs}
Les premiers fluides de réfrigérateurs, comme le dichlorodifluorométhane $\ce{CCl2F2}$, détruisaient la couche d'ozone : ils ont été
interdits par le protocole de Montréal (1987). On les a remplacés par des fluides sans chlore, comme le difluorométhane
$\ce{CH2F2}$. Certaines installations utilisent aujourd'hui le dioxyde de carbone.
\end{docu}
\begin{docu}{2}{Électronégativités}
\par\smallskip\noindent\hfill\begin{tabular}{lcccc}
\toprule
Élément & H & C & Cl & F \\
\midrule
$\chi$ & 2,20 & 2,55 & 3,16 & 3,98 \\
\bottomrule
\end{tabular}\hfill\null\par\smallskip
\end{docu}
\begin{questions}
\item Calculer le nombre de doublets de $\ce{CCl2F2}$ et de $\ce{CH2F2}$, puis établir leurs schémas de Lewis. L'atome central est le carbone.
\item Prévoir leur géométrie.
\item Identifier les liaisons polarisées dans chaque molécule.
\item \initiative Parmi les trois fluides cités, lequel est apolaire ? Justifier.
\end{questions}
\end{exercice}
\end{blocLivre}
\begin{blocLivre}{ex-6-14}
\begin{exercice}[3]{Un gaz incolore et mortel}
\begin{docu}{1}{Le monoxyde de carbone}
Le monoxyde de carbone $\ce{CO}$ se forme lors des combustions incomplètes. Incolore et inodore, il se fixe sur l'hémoglobine du
sang à la place du dioxygène : il cause chaque année plusieurs dizaines de décès en France.
\end{docu}
\begin{docu}{2}{Deux molécules jumelles}
Le diazote $\ce{N2}$ et le monoxyde de carbone possèdent le même nombre d'électrons de valence. Dans le schéma de Lewis de $\ce{N2}$,
chaque atome d'azote porte un doublet non liant.
Une mesure montre que, dans $\ce{CO}$, les centres des charges positives et négatives ne sont pas confondus.
\end{docu}
\begin{questions}
\item Établir le schéma de Lewis du diazote et vérifier la règle de l'octet.
\item Vérifier que $\ce{CO}$ et $\ce{N2}$ ont le même nombre d'électrons de valence.
\item \initiative Proposer un schéma de Lewis pour le monoxyde de carbone, chaque atome respectant la règle de l'octet. Quelle
particularité présente l'atome de carbone ?
\item Quelle est la géométrie de la molécule ? Est-elle polaire ?
\end{questions}
\end{exercice}
\end{blocLivre}
\begin{blocLivre}{ex-6-15}
\begin{exercice}[3]{Identifier des modèles moléculaires}
\begin{docu}{1}{Trois modèles}
Un élève a construit trois molécules avec une boîte de modèles.
\schema{06-ex15}
\end{docu}
\begin{questions}
\item Écrire la formule brute de chaque molécule.
\item \initiative Établir leurs schémas de Lewis et vérifier qu'ils sont compatibles avec les modèles.
\item Nommer la géométrie de chaque molécule, ou autour de chaque atome de carbone.
\item Quelles molécules sont polaires ?
\end{questions}
\end{exercice}
\end{blocLivre}
